\section*{Chapitre \ref{ch:b1:hotspots} --- Points chauds et volcanisme intraplaque}

\begin{solution}{exo:b1:hotspots:1}
(a) $4205 + 5274 = \qty{9479}{m}$, près de \qty{9.5}{km}. Le volcan est
plus lourd que l'eau qu'il remplace~: la plaque fléchit sous sa charge, donc
sa base se trouve plus profond que le plancher lointain, et sa vraie hauteur est plus grande.
(b) $74000/0.1 = \num{740000}$ ans à l'apport du Kilauea~;
$74000/0.025 = \num{3.0e6}$ ans au rythme moyen de la chaîne. (c)
$0.12\times19 = \qty{2.3}{km^3}$, environ $\num{3e-5}$ du Mauna Loa~: trois
millièmes de un pour cent.
\end{solution}

\begin{solution}{exo:b1:hotspots:2}
(a) $0.92\times\pi/180 = \qty{0.0161}{rad.\Ma^{-1}}$. (b) En radians par
million d'années fois kilomètres, $\omega R\sin\Delta$ est en kilomètres par
million d'années, c'est-à-dire en millimètres par an~: $0.016057\times6371 = 102.3$, donc
$v = 102.3\sin\Delta$, ce qui donne \qty{51.1}{mm.an^{-1}} à \ang{30},
\qty{88.6}{mm.an^{-1}} à \ang{60}, \qty{102.3}{mm.an^{-1}} à \ang{90} et
\qty{97.3}{mm.an^{-1}} à \ang{108.0}. (c) La vitesse est proportionnelle à la
distance $R\sin\Delta$ à l'axe, maximale sur le \omterm{def:b1:spherical-astronomy:spherical-triangle}{grand cercle} à \ang{90}
du pôle~; au-delà, le point se rapproche de nouveau de l'axe, du
côté du pôle opposé.
\end{solution}

\begin{solution}{exo:b1:hotspots:3}
(a) Deccan~: $\num{e6}/\num{5e5} = \qty{2}{km}$. Atlantique central~:
$\num{2e6}/\num{7e6} = \qty{0.29}{km}$, environ \qty{290}{m}. (b) La moyenne pour
le Deccan est proche de sa coupe de \qty{2.5}{km}~: la province est un plateau épais et
continu. (c) L'aire de l'Atlantique central
est celle de toute la région, sur quatre continents, où l'on trouve ses
restes~: laves, dykes et filons-couches maintenant séparés par un océan et très
érodés. La roche est concentrée en certains endroits et absente de la plus grande
part de l'aire, donc la moyenne ne dit rien d'un empilement isolé.
\end{solution}

\begin{solution}{exo:b1:hotspots:4}
(a) $\rho_m\alpha\,\Delta T = 3300\times\num{3.2e-5}\times250 =
\qty{26.4}{kg.m^{-3}}$. (b) La colonne du renflement a la masse supplémentaire
$(\rho_m - \rho_w)h = 2300\times690 = \qty{1.587e6}{kg.m^{-2}}$, que la
couche légère doit équilibrer~: $H = \num{1.587e6}/26.4 = \qty{6.0e4}{m}$, environ
\qty{60}{km} de roche chaude.
\end{solution}

\begin{solution}{exo:b1:hotspots:5}
(a) En A ($\lambda = 0$, $\varphi = 0$), $\vec e = (0, 1, 0)$ et
$\vec n = (0, 0, 1)$~; avec $\sin\ang{270} = -1$ et $\cos\ang{270} = 0$,
$\vec t_A = -\vec e = (0, -1, 0)$. En B ($\lambda = \ang{90}$),
$\vec e = (-1, 0, 0)$, $\vec n = (0, 0, 1)$ et $\vec t_B = (1, 0, 0)$.
(b) $\vec t_A\times\vec t_B = (0, 0, 1)$. (c) Avec $\vec p' = (0, 0, 1)$ et
$\vec r_A = (1, 0, 0)$, $\vec p'\times\vec r_A = (0, 1, 0)$, et
$(\vec p'\times\vec r_A)\cdot\vec t_A = -1 < 0$~: le pôle est
$-\vec p' = (0, 0, -1)$, le pôle Sud. Un mouvement vers l'ouest le long de l'équateur
est horaire vu de dessus le pôle Nord, donc antihoraire vu de dessus
le pôle Sud. (d) À \ang{90} du pôle, $\omega = 50/6371 =
\qty{7.85e-3}{rad.\Ma^{-1}} = \qty{0.45}{\degree\per\Ma}$.
\end{solution}

\begin{solution}{exo:b1:hotspots:6}
(a) $\sum x_i v_i = 90000 + 240000 + 342000 = \num{672000}$ et
$\sum x_i^2 = \num{7.0e7}$, donc $\omega = \qty{0.0096}{rad.\Ma^{-1}}$, soit
\qty{0.55}{\degree\per\Ma}. (b) Les prédictions sont 28.8, 48.0 et
\qty{57.6}{mm.an^{-1}}~; les résidus $+1.2$, $0.0$ et $-0.6$
\unit{mm.an^{-1}}. (c) Pour une rotation rigide la vitesse est proportionnelle à la
distance à l'axe, et un point sur l'axe ne bouge pas~: il n'y a pas
d'ordonnée à l'origine à ajuster.
\end{solution}

\begin{solution}{exo:b1:hotspots:7}
(a) $v = \qty{50e-3}{m}/\qty{3.156e7}{s} = \qty{1.58e-9}{m.s^{-1}}$, donc
$B = 2300\times\num{3e8}\times\num{1.58e-9} = \qty{1090}{kg.s^{-1}}$ et
$Q = (910/\num{3.2e-5})\times1090 = \qty{3.1e10}{W}$. (b) Un tiers du
flux hawaïen. (c) $B$ croît comme $A\,v$~: le panache sous la plaque plus rapide
doit refaire le même renflement deux fois plus vite, et son flux est deux fois plus grand.
\end{solution}

\begin{solution}{exo:b1:hotspots:8}
(a) $350/(16.1 - 10.3) = 350/5.8 = \qty{60}{km.\Ma^{-1}}$, c'est-à-dire
\qty{60}{mm.an^{-1}}. (b) La jeune trace, \qty{29 \pm 5}{mm.an^{-1}}, concorde
avec les \qty{26 \pm 5}{mm.an^{-1}} de la plaque dans les incertitudes~; la
partie ancienne est allée plus de deux fois plus vite, $60 - 26 = 34$, presque sept incertitudes-types
de la vitesse de la plaque. (c) Oui~: N\,\ang{54}\,E est l'\omterm{def:b1:spherical-astronomy:horizontal}{azimut}
\ang{54}, opposé à \ang{235} à un degré près. Le volcanisme se déplace sur la
plaque dans le sens opposé à celui de la plaque sur le \omterm{def:b1:hotspots:hotspot}{point chaud}. (d) L'étirement a écarté
les caldeiras plus anciennes, allongeant les distances entre elles et donc la vitesse
apparente~; la plaque elle-même n'est pas allée plus vite.
\end{solution}

\begin{solution}{exo:b1:hotspots:9}
(a) $\num{9.2e10}/\num{e12} = \qty{0.092}{W.m^{-2}}$, \qty{92}{mW.m^{-2}}.
(b) Neuf à dix-huit fois l'excès observé, qui peut même être nul.
(c) $L^2/\kappa =
\num{e10}/\num{e-6} = \qty{e16}{s}$, environ \qty{320}{\Ma}~; la plaque parcourt
\qty{1000}{km} en $1000/92 = \qty{11}{\Ma}$. La chaleur venue d'en dessous n'a eu le temps de
traverser qu'une petite partie de la plaque pendant qu'elle est au-dessus du renflement~: l'essentiel
de la chaleur du panache reste dans le \omterm{def:b1:layered-earth:layers}{manteau} sous la plaque et est emporté avec
elle, et la surface en voit peu.
\end{solution}

\begin{solution}{exo:b1:hotspots:10}
(a) Dix fois le liquide fondu, \qty{e7}{km^3}, de rayon
$(3V/4\pi)^{1/3} = \qty{134}{km}$. (b) $\num{e7}/100 = \qty{e5}{km^2}$, un disque
de $2\sqrt{\num{e5}/\pi} = \qty{357}{km}$ de large~: dix-huit fois plus petit en aire
que les trapps. La tête était plus grande ou fondait moins, ou elle s'étalait plus large et
plus mince, ou les laves ont coulé bien au-delà d'elle. (c) \qty{1}{km^3} par an, dix
fois l'apport actuel du Kilauea. (d) Un nouveau panache doit rassembler assez de
roche chaude et légère pour se frayer un chemin à travers le \omterm{def:b1:layered-earth:layers}{manteau} visqueux, et un grand
volume monte bien plus vite qu'un petit~; la queue qui suit remonte par le
chemin chaud et faible que la tête a ouvert, et un conduit étroit suffit.
\end{solution}

\begin{solution}{exo:b1:hotspots:11}
(a) $\qty{40}{km.\Ma^{-1}}\times(81 - 47)\,\unit{\Ma} = \qty{1360}{km}$.
(b) La plaque se déplace à $(-80\sin\ang{20}, 80\cos\ang{20}) = (-27.4, 75.2)$
et le \omterm{def:b1:hotspots:hotspot}{point chaud} à $(0, -40)$~: par rapport au \omterm{def:b1:hotspots:hotspot}{point chaud} la plaque se déplace à
$(-27.4, 115.2)$ \unit{mm.an^{-1}}, une vitesse de \qty{118.4}{mm.an^{-1}} vers
$\ang{360} - \arctan(27.4/115.2) = \ang{346.6}$. (c) La trace donnerait une
vitesse trop élevée de \qty{48}{\percent} et une direction trop à l'est de \ang{6.6}~:
un \omterm{def:b1:hotspots:hotspot}{point chaud} en dérive enregistre le mouvement de la plaque par rapport à lui-même, non par rapport au
\omterm{def:b1:layered-earth:layers}{manteau} profond.
\end{solution}

\begin{solution}{exo:b1:hotspots:12}
(a) $\qty{3e6}{J.m^{-3}.K^{-1}}\times\qty{1300}{K}\times\qty{e15}{m^3} =
\qty{3.9e24}{J}$. (b) Sur $\qty{3.16e13}{s}$, \qty{1.2e11}{W}~: environ la
chaleur portée par le panache hawaïen aujourd'hui (1.3 fois), et \qty{0.26}{\percent}
des \qty{47}{TW} de la Terre. (c) Une borne inférieure~: la tête apporte aussi
la chaleur excédentaire des neuf dixièmes de sa roche qui ne fondent pas, et réchauffe la
base de la \omterm{def:b1:layered-earth:lithosphere}{lithosphère}, et la chaleur latente a été laissée de côté.
\end{solution}

\begin{solution}{pb:b1:hotspots:1}
\textbf{1.} À Hawaï ($\varphi = \ang{19.0}$, $\lambda = \ang{-155.2}$),
avec $\sin\ang{304} = -0.8290$ et $\cos\ang{304} = 0.5592$~:
\[
\begin{aligned}
  \vec r_H &= (-0.8583,\ -0.3966,\ 0.3256), &
  \vec e &= (0.4195,\ -0.9078,\ 0),\\
  \vec n &= (0.2955,\ 0.1366,\ 0.9455), &
  \vec t_H &= (-0.1825,\ 0.8289,\ 0.5287).
\end{aligned}
\]

\textbf{2.} À Louisville ($\varphi = \ang{-53.6}$, $\lambda = \ang{-140.6}$),
avec $\sin\ang{316} = -0.6947$ et $\cos\ang{316} = 0.7193$~:
\[
\begin{aligned}
  \vec r_L &= (-0.4586,\ -0.3767,\ -0.8049), &
  \vec e &= (0.6347,\ -0.7727,\ 0),\\
  \vec n &= (-0.6220,\ -0.5109,\ 0.5934), &
  \vec t_L &= (-0.8883,\ 0.1693,\ 0.4269).
\end{aligned}
\]

\textbf{3.} $\lVert\vec t_H\rVert^2 = 0.0333 + 0.6871 + 0.2795 = 0.9999$ et
$\vec t_H\cdot\vec r_H = 0.1566 - 0.3287 + 0.1721 = 0.0000$, aux arrondis près.

\textbf{4.} Le produit vectoriel est $\vec t_H\times\vec t_L = (0.2643, -0.3917, 0.7054)$, de longueur
$0.8491$. L'angle entre les traces est $\arcsin0.8491 = \ang{58.1}$ (leur
produit scalaire, $0.528$, est positif, donc l'angle est aigu).

\textbf{5.} Le pôle provisoire est $\vec p' = (0.3113, -0.4613, 0.8308)$. Alors
$(\vec p'\times\vec r_H)\cdot\vec t_H = -0.98 < 0$~: le pôle est
$\vec p = (-0.3113, 0.4613, -0.8308)$.

\textbf{6.} $\arcsin(-0.8308) = \ang{-56.18}$ et
$\operatorname{atan2}(0.4613, -0.3113) = \ang{124.01}$~: le pôle se trouve à
\ang{56.2}\,S, \ang{124.0}\,E, dans l'océan Austral au sud de l'Australie.

\textbf{7.} $\vec p\cdot\vec r_H = 0.2672 - 0.1830 - 0.2705 = -0.1863$, donc
$\Delta_H = \ang{100.7}$~; $\vec p\cdot\vec r_L = 0.1428 - 0.1738 + 0.6687 =
0.6377$, donc $\Delta_L = \ang{50.4}$.

\textbf{8.} Le vecteur $\vec p\times\vec r_L$ est perpendiculaire à $\vec p$ et
à $\vec r_L$, comme $\vec t_L$ (par construction de $\vec p$, et parce qu'une
trace est horizontale)~; les deux sont le long de la même droite, et sa longueur est
$\sin\Delta_L$. Donc $(\vec p\times\vec r_L)\cdot\vec t_L = \pm\sin\Delta_L$, et
le calcul donne $+0.770 = \sin\ang{50.4}$~: la plaque se déplace à Louisville
selon $+\vec t_L$, dans le sens où va la trace. La règle de signe, appliquée à Hawaï,
aurait pu donner le mauvais sens à Louisville~; il n'en est rien, donc les deux
directions sont reproduites exactement, \ang{304} et \ang{316}.

\textbf{9.} $\omega = 67/(6371\sin\ang{50.38}) = 67/4908 =
\qty{0.01365}{rad.\Ma^{-1}}$, c'est-à-dire \qty{0.782}{\degree\per\Ma}.

\textbf{10.} $360/0.782 = \qty{460}{\Ma}$.

\textbf{11.} La formule des cosinus entre \ang{56.18}\,S, \ang{124.01}\,E et
\ang{63.78}\,S, \ang{112.48}\,E donne \ang{9.5}. La vitesse à deux traces,
\qty{0.78}{\degree\per\Ma}, est inférieure de \qty{15}{\percent} aux
\qty{0.92}{\degree\per\Ma} ajustés sur treize vitesses.

\textbf{12.} Société \ang{73.6}, Pitcairn \ang{77.8}, Cobb \ang{134.7}.

\textbf{13.} On calcule $\omega R = 0.01365\times6371 = \qty{86.98}{mm.\year^{-1}}$, fois
$\sin\Delta$~: Société \qty{83.5}{mm.an^{-1}}, Pitcairn
\qty{85.0}{mm.an^{-1}}, Cobb \qty{61.8}{mm.an^{-1}}.

\textbf{14.} Pitcairn concorde~: \qty{5}{mm.an^{-1}} de moins que la mesure, dans
ses \qty{15}{mm.an^{-1}}. La Société est trop lente de \qty{25.5}{mm.an^{-1}}, 2.6
incertitudes, et Cobb trop rapide de \qty{18.8}{mm.an^{-1}}, plus de six~:
ni l'une ni l'autre ne concorde.

\textbf{15.} Le pôle est perpendiculaire aux deux vecteurs de trace. S'ils sont
presque parallèles, leur produit vectoriel est court, et une erreur de quelques degrés sur
l'une ou l'autre direction le fait tourner d'un grand angle~: le pôle est mal
déterminé.

\textbf{16.} $v = \qty{92e-3}{m}/\qty{3.156e7}{s} = \qty{2.92e-9}{m.s^{-1}}$
et $B = 2300\times\num{4.83e8}\times\num{2.92e-9} = \qty{3240}{kg.s^{-1}}$.

\textbf{17.} $Q = (910/\num{3.2e-5})\times3240 = \qty{9.2e10}{W}$,
\qty{0.20}{\percent} de \qty{47}{TW}.

\textbf{18.} Les dorsales font $21/67000 = \qty{3.1e-4}{km^3}$ par an par
kilomètre~; $0.1/\num{3.1e-4} = \qty{320}{km}$ de dorsale.

\textbf{19.} $74000/0.1 = \num{740000}$ ans.

\textbf{20.} $\num{2e6}/\num{81e6} = \qty{0.025}{km^3}$ par an~: le Kilauea seul
est alimenté aujourd'hui avec quatre fois la production moyenne de tout le \omterm{def:b1:hotspots:hotspot}{point chaud}, qui
est près de son maximum.

\textbf{21.} $v_H = 86.98\sin\ang{100.73} = 86.98\times0.9825 =
\qty{85.5}{mm.\year^{-1}}$.

\textbf{22.} $92 - 85.5 = \qty{6.5}{mm.\year^{-1}}$, 2.2 incertitudes-types de
la vitesse mesurée~: les deux \omterm{def:b1:hotspots:hotspot}{points chauds} ne s'accordent pas tout à fait sur le mouvement d'une
\omterm{def:b1:plate-kinematics:rigid-plate}{plaque rigide} sur un repère fixe.

\textbf{23.} Les âges donnent \qty{9.1}{cm} et \qty{8.6}{cm} par an, moyennés
sur toute la chaîne~; les \qty{9.2}{cm} par an de Morgan et Morgan sont proches
de la première, et la prédiction de Louisville, \qty{8.5}{cm} par an, de la
seconde. Les quatre sont à moins de \qty{0.7}{cm} par an les unes des autres~: traces, âges
et directions racontent la même histoire à environ \qty{8}{\percent} près.

\textbf{24.} À une décimale, la trace de Louisville donne à la plaque pacifique au-dessus
du \omterm{def:b1:hotspots:hotspot}{point chaud} hawaïen la vitesse
\textbf{\qty[text-series-to-math=true]{8.5}{cm.an^{-1}}}.
\end{solution}
